% 平面仿射半空间 h_{w,b}(x)=I{<w,x>+b>=0} 的打散与不可实现标记（第5章5.2.3）。
% 左侧：不共线三点 P1=(5,3)、P2=(15,3)、P3=(9,12)（各 20x16 小格的局部坐标，单位 mm）上的全部 8 种标记，
%   每格画出一条实现该标记的分界线，浅蓝区域为预测 1 的半平面；全 0 与全 1 用整体偏移的半平面实现。
%   分界线：{P1}: 过 (9,0) 与 (0,9)；{P2}: 过 (11,0) 与 (20,9)；{P3}: y=8；{P1,P2}: y=7；{P1,P3}: 过 (12,0) 与 (16,16) 的直线左侧；{P2,P3}: 过 (10,0) 与 (3,16) 的直线右侧。
% 右侧：四点的两种几何情形。(a) D=(10,6) 在 A=(0,0)、B=(20,0)、C=(10,16) 的凸包内，标记 A,B,C 为 1、D 为 0 不可实现；
%   (b) 凸四边形 E=(0,0)、F=(20,4)、G=(16,16)、K=(2,12)，对角线 EG（y=x）与 FK 交于 (116/13,116/13)≈(8.92,8.92)，
%   标记 E,G 为 1、F,K 为 0 不可实现。实心蓝点为 1，空心红圈为 0。
\begingroup
\begin{tikzpicture}[figure picture,x=1mm,y=1mm,
  one/.style={circle,fill=FigureBlue,draw=FigureBlue,inner sep=0pt,minimum size=2.3mm},
  zero/.style={circle,fill=white,draw=FigureRed,line width=.75pt,inner sep=0pt,minimum size=2.3mm},
  cut/.style={draw=FigureStroke,line width=.75pt},
  hp positive/.style={fill=FigureBlueFill}]
\node[figure heading,anchor=west] at (0,47) {不共线的三点：8 种标记都可实现};
% 小格辅助命令：#1,#2 为小格原点，#3 填色，#4 分界线，#5 三个点，#6 小格说明。
\newcommand{\hpPanel}[6]{% #1 ox, #2 oy, #3 fill path (relative), #4 cut path, #5 labels P1..P3 styles as a/b/c, #6 caption
  \begin{scope}[shift={(#1,#2)}]
    \begin{scope}\clip (0,0) rectangle (20,16); #3 \end{scope}
    \draw[FigureGrid,line width=.5pt] (0,0) rectangle (20,16);
    #4
    #5
    \node[figure note] at (10,-3.2) {#6};
  \end{scope}}
% points in panel coords: P1=(5,3) P2=(15,3) P3=(9,12)
\def\Pa{(5,3)}\def\Pb{(15,3)}\def\Pc{(9,12)}
\hpPanel{0}{24}{}{}{\node[zero] at \Pa{};\node[zero] at \Pb{};\node[zero] at \Pc{};}{全为 0}
\hpPanel{24}{24}{\fill[hp positive] (0,0)--(9,0)--(0,9)--cycle;}{\draw[cut] (9,0)--(0,9);}{\node[one] at \Pa{};\node[zero] at \Pb{};\node[zero] at \Pc{};}{$P_1$ 为 1}
\hpPanel{48}{24}{\fill[hp positive] (11,0)--(20,0)--(20,9)--cycle;}{\draw[cut] (11,0)--(20,9);}{\node[zero] at \Pa{};\node[one] at \Pb{};\node[zero] at \Pc{};}{$P_2$ 为 1}
\hpPanel{72}{24}{\fill[hp positive] (0,8)rectangle(20,16);}{\draw[cut] (0,8)--(20,8);}{\node[zero] at \Pa{};\node[zero] at \Pb{};\node[one] at \Pc{};}{$P_3$ 为 1}
\hpPanel{0}{0}{\fill[hp positive] (0,0)rectangle(20,7);}{\draw[cut] (0,7)--(20,7);}{\node[one] at \Pa{};\node[one] at \Pb{};\node[zero] at \Pc{};}{$P_1,P_2$ 为 1}
\hpPanel{24}{0}{\fill[hp positive] (0,0)--(12,0)--(16,16)--(0,16)--cycle;}{\draw[cut] (12,0)--(16,16);}{\node[one] at \Pa{};\node[zero] at \Pb{};\node[one] at \Pc{};}{$P_1,P_3$ 为 1}
\hpPanel{48}{0}{\fill[hp positive] (10,0)--(20,0)--(20,16)--(3,16)--cycle;}{\draw[cut] (10,0)--(3,16);}{\node[zero] at \Pa{};\node[one] at \Pb{};\node[one] at \Pc{};}{$P_2,P_3$ 为 1}
\hpPanel{72}{0}{\fill[hp positive] (0,0)rectangle(20,16);}{}{\node[one] at \Pa{};\node[one] at \Pb{};\node[one] at \Pc{};}{全为 1}
\node[figure note,anchor=west] at (0,42.5) {$P_1$ 左下，$P_2$ 右下，$P_3$ 居上；浅蓝区域为预测 1 的半平面};
% right part
\node[figure heading,anchor=west] at (102,47) {任意四点：总有一种标记不可实现};
\begin{scope}[shift={(104,14)}]
  \fill[FigureBlueFill] (0,0)--(20,0)--(10,16)--cycle;
  \draw[FigureBlue,dash pattern=on 2pt off 1.5pt] (0,0)--(20,0)--(10,16)--cycle;
  \node[one] at (0,0) {};\node[one] at (20,0) {};\node[one] at (10,16) {};\node[zero] at (10,6) {};
  \node[figure note,align=center] at (10,-7.5) {(a) 一点在其余三点\\的凸包内};
\end{scope}
\begin{scope}[shift={(136,14)}]
  \draw[FigureGrid,line width=.6pt] (0,0)--(20,4)--(16,16)--(2,12)--cycle;
  \draw[FigureBlue,line width=.9pt] (0,0)--(16,16);
  \draw[FigureRed,line width=.9pt,dash pattern=on 2pt off 1.5pt] (20,4)--(2,12);
  \fill[FigureInk] (8.92,8.92) circle (.6);
  \node[one] at (0,0) {};\node[one] at (16,16) {};\node[zero] at (20,4) {};\node[zero] at (2,12) {};
  \node[figure note,align=center] at (10,-7.5) {(b) 凸四边形，对角线\\两端分别同标};
\end{scope}
\end{tikzpicture}
\endgroup
